\documentclass[aspectratio=169,UTF8]{ctexbeamer}
\usetheme{Madrid}
\usecolortheme{default}
\definecolor{DeepBlue}{HTML}{2563EB}
\definecolor{DeepInk}{HTML}{172033}
\definecolor{DeepGreen}{HTML}{059669}
\setbeamercolor{structure}{fg=DeepBlue}
\setbeamercolor{title}{fg=white,bg=DeepBlue}
\setbeamercolor{frametitle}{fg=white,bg=DeepBlue}
\setbeamercolor{alerted text}{fg=DeepGreen}
\setbeamertemplate{navigation symbols}{}
\setbeamertemplate{footline}[frame number]
\usepackage{booktabs}
\usepackage{listings}
\usepackage{xcolor}
\usepackage{hyperref}
\hypersetup{colorlinks=true,urlcolor=DeepBlue}
\lstset{
  basicstyle=\ttfamily\scriptsize,
  keywordstyle=\color{DeepBlue},
  stringstyle=\color{DeepGreen},
  commentstyle=\color{gray},
  showstringspaces=false,
  breaklines=true,
  frame=single,
  columns=fullflexible
}
\newcommand{\api}[1]{\texttt{#1}}
\newcommand{\keypoint}[1]{\alert{#1}}
\AtBeginSection[]{
  \begin{frame}{本节路线}
    \tableofcontents[currentsection]
  \end{frame}
}

\title{第5课：结构化输出与补全能力}
\subtitle{JSON Output、对话前缀续写与 FIM}
\author{DeepSeek API 全栈课程}
\date{文档快照：2026-11-05}
\begin{document}
\frame{\titlepage}
\begin{frame}{学习目标}
\begin{itemize}
  \item 生成可解析 JSON 并处理空白/截断
  \item 用 assistant 前缀约束回答开头
  \item 用 FIM 补全代码中间部分
  \item 认识 logprobs、Beta 地址和能力边界
\end{itemize}
\end{frame}
\section{JSON Output}
\begin{frame}{开启 JSON 模式}
设置 \api{response\_format=\{"type":"json\_object"\}}，并在 system 或 user 提示中明确包含“JSON”字样及目标结构示例。\\[8pt]
API 保证输出是合法 JSON 字符串，但\keypoint{不保证字段满足业务语义}；仍需客户端 Schema 校验。
\end{frame}
\begin{frame}[fragile]{结构化请求}
\begin{lstlisting}[language=Python]
response = client.chat.completions.create(
  model="deepseek-flash",
  messages=[
    {"role": "system", "content":
      "输出 JSON：{\"name\": str, \"score\": int}"},
    {"role": "user", "content": "小明得 95 分"},
  ],
  response_format={"type": "json_object"},
  max_tokens=200,
)
data = json.loads(response.choices[0].message.content)
\end{lstlisting}
\end{frame}
\begin{frame}{四层校验}
\begin{enumerate}
  \item \api{finish\_reason != length}
  \item content 非空且不是纯空格
  \item \api{json.loads} 成功
  \item 业务 Schema：必填字段、类型、范围、枚举
\end{enumerate}
失败时缩短输入/扩大预算、改进示例后有限重试；不要无限重试同一提示。
\end{frame}
\begin{frame}{特殊空白行为}
若模型输出无法识别为 JSON，API 会把非 JSON 内容替换为 U+0020 空格；content 看似为空，但 finish reason 仍可能为 stop。\\[8pt]
因此 \api{if not content.strip()} 是必需的防线。
\end{frame}
\section{对话前缀续写}
\begin{frame}{Prefix Completion（Beta）}
\begin{itemize}
  \item 沿用 Chat Completions
  \item base URL 改为 \api{https://api.deepseek.com/beta}
  \item 最后一条必须是 assistant 消息
  \item 该消息设置 \api{prefix=True}
\end{itemize}
用于强制代码围栏、固定措辞开头或模板续写。
\end{frame}
\begin{frame}[fragile]{前缀 + stop}
\begin{lstlisting}[language=Python]
messages = [
  {"role": "user", "content": "写快速排序"},
  {"role": "assistant", "content": "```python\n",
   "prefix": True},
]
response = client.chat.completions.create(
  model="deepseek-flash", messages=messages,
  stop=["```"],
)
\end{lstlisting}
前缀只是生成约束，不是代码安全沙箱。
\end{frame}
\section{FIM 补全}
\begin{frame}{Fill In the Middle（Beta）}
FIM 使用 \api{POST /beta/completions}：\api{prompt} 是前缀，\api{suffix} 是后缀，模型生成中间内容。适合 IDE 补全和局部模板填充。\\[8pt]
最大补全长度为 4K；仅非思考模式支持。
\end{frame}
\begin{frame}{Completion 参数与响应}
\begin{itemize}
  \item 输入：model、prompt、suffix、max\_tokens、stop
  \item 可选：echo、logprobs、temperature、top\_p、stream
  \item 输出文本：\api{choices[0].text}
  \item 用量结构与 Chat 类似，含缓存明细
\end{itemize}
\end{frame}
\begin{frame}{FIM 上下文设计}
\begin{itemize}
  \item 前后缀应保留函数签名、缩进、类型与邻近依赖
  \item stop 阻止越界生成到下一函数/围栏
  \item 生成代码必须格式化、静态检查、测试后再采用
  \item 不把仓库密钥或无关私有代码发送给模型
\end{itemize}
\end{frame}
\section{概率与选择}
\begin{frame}{什么时候用哪一种}
\begin{tabular}{ll}
\toprule
目标 & 能力\\
\midrule
可靠机器读取 & JSON Output + 本地 Schema 校验\\
固定回答开头 & Chat Prefix（Beta）\\
补全中间代码 & FIM（Beta）\\
观察候选 Token & logprobs / top\_logprobs\\
\bottomrule
\end{tabular}
\end{frame}
\begin{frame}{本课例程}
\begin{enumerate}
  \item \api{01\_json\_output.py}：解析与兜底
  \item \api{02\_prefix\_completion.py}：代码前缀续写
  \item \api{03\_fim\_completion.py}：中间补全
\end{enumerate}
\end{frame}
\begin{frame}{官方资料}
\small
\url{https://api-docs.deepseek.com/zh-cn/guides/json_mode}\\
\url{https://api-docs.deepseek.com/zh-cn/guides/chat_prefix_completion}\\
\url{https://api-docs.deepseek.com/zh-cn/guides/fim_completion}\\
\url{https://api-docs.deepseek.com/zh-cn/api/create-completion}
\end{frame}
\end{document}
